% Experiment Template
% 1) Copy this file and rename it, e.g. experiments/experiment_04_firewall_rules.tex
% 2) In main.tex, add:
%    \includeexperiment{XX}{Experiment Title}{experiments/your_new_file_name}
% 3) Keep this file content as sections only (no \experiment command here)

\section{Theory}
The Vigen\`ere cipher is a polyalphabetic substitution cipher that uses a string key converted to a list of integers (0--25). Each plaintext character is shifted by the corresponding key integer using modular addition.

Encryption: $C_i = (P_i + K_{i \bmod m}) \bmod 26$, where $m$ is the key length. Decryption reverses this: $P_i = (C_i - K_{i \bmod m} + 26) \bmod 26$. For manual computation, a Vigen\`ere table (Tabula Recta) can be used, but the program computes the cipher directly using modular arithmetic.

\section{Algorithm}
Encryption Algorithm of Vigenere Cipher is,
\begin{enumerate}
  \item set mod := length(key)
  \item set len := length(plaintext)
  \item for i := 0 to len-1 do
  \begin{enumerate}
    \item if plaintext[i] is an uppercase letter then
    \begin{enumerate}
      \item cipher[i] := ((plaintext[i] - 'A') + key[i \% mod]) mod 26 + 'A'
    \end{enumerate}
    \item else if plaintext[i] is a lowercase letter then
    \begin{enumerate}
      \item cipher[i] := ((plaintext[i] - 'a') + key[i \% mod]) mod 26 + 'a'
    \end{enumerate}
    \item else
    \begin{enumerate}
      \item cipher[i] := plaintext[i]
    \end{enumerate}
  \end{enumerate}
  \item return cipher;
\end{enumerate}

Decryption Algorithm of Vigenere Cipher is,
\begin{enumerate}
  \item set mod := length(key)
  \item set len := length(ciphertext)
  \item for i := 0 to len-1 do
  \begin{enumerate}
    \item if ciphertext[i] is an uppercase letter then
    \begin{enumerate}
      \item plaintext[i] := ((ciphertext[i] - 'A') - key[i \% mod] + 26) mod 26 + 'A'
    \end{enumerate}
    \item else if ciphertext[i] is a lowercase letter then
    \begin{enumerate}
      \item plaintext[i] := ((ciphertext[i] - 'a') - key[i \% mod] + 26) mod 26 + 'a'
    \end{enumerate}
    \item else
    \begin{enumerate}
      \item plaintext[i] := ciphertext[i]
    \end{enumerate}
  \end{enumerate}
  \item return plaintext;
\end{enumerate}

% \section{Flowchart}
% \begin{figure}[H]
%   \centering
%   \begin{autofitflowchart}
%     % ---------- Encryption ----------
%     \node (es)   [terminator]                         {Start};
%     \node (ein)  [io,      right=2.2cm of es]         {Input plaintext, key};
%     \node (ekey) [process, right=2.2cm of ein]        {Build key[],\; mod $\leftarrow$ length(key)};

%     \node (einit) [process, below=1.7cm of es]        {$i \leftarrow 0,\; len \leftarrow$ length(plaintext)};
%     \node (eloop) [decision,below=1.7cm of ein]       {$i < len$ ?};
%     \node (eout)  [io,      below=1.7cm of ekey]      {Output cipher};

%     \node (eupq)  [decision,below=1.8cm of ein,xshift=-2.2cm] {Uppercase?};
%     \node (elowq) [decision,below=1.8cm of ein]              {Lowercase?};
%     \node (ecopy) [process, below=1.8cm of ein,xshift= 2.2cm] {cipher[i] $\leftarrow$ plaintext[i]};

%     \node (eup)   [process, below=1.8cm of eupq]      {cipher[i] $\leftarrow$ ((p[i]-`A') + key[i \% mod]) \% 26 + `A'};
%     \node (elow)  [process, below=1.8cm of elowq]     {cipher[i] $\leftarrow$ ((p[i]-`a') + key[i \% mod]) \% 26 + `a'};
%     \node (einc)  [process, below=1.8cm of ecopy]     {$i \leftarrow i+1$};

%     \node (eret)  [process, below=1.8cm of eup]       {Return cipher};
%     \node (eend)  [terminator, right=2.2cm of eret]   {End};
%     \node (epad)  [process, right=2.2cm of eend]      {Encryption complete};

%     % Main flow
%     \draw [line] (es) -- (ein);
%     \draw [line] (ein) -- (ekey);
%     \draw [line] (ekey) |- (einit);
%     \draw [line] (einit) -- (eloop);

%     % Loop decision
%     \draw [line] (eloop.east) -- node[above]{No} (eout.west);
%     \draw [line] (eout) |- (eret);

%     \draw [line] (eloop.south) -- node[right]{Yes} (elowq.north);
%     \draw [line] (elowq.west) -- node[above]{No} (eupq.east);
%     \draw [line] (elowq.east) -- node[above]{No} (ecopy.west);
%     \draw [line] (elowq.south) -- node[right]{Yes} (elow.north);

%     \draw [line] (eupq.south) -- node[right]{Yes} (eup.north);

%     \draw [line] (eup.east) -- ++(0.9,0) |- (einc.west);
%     \draw [line] (elow.east) -- ++(0.9,0) |- (einc.west);
%     \draw [line] (ecopy.south) -- (einc.north);

%     % Loop back (routed outside, no node crossing)
%     \draw [line] (einc.east) -- ++(1.2,0) |- (eloop.east);

%     % End path
%     \draw [line] (eret) -- (eend);
%     \draw [line] (eend) -- (epad);
%   \end{autofitflowchart}
%   \caption{Vigenere Cipher Encryption Flowchart}
% \end{figure}

% \begin{figure}[H]
%   \centering
%   \begin{autofitflowchart}
%     % ---------- Decryption ----------
%     \node (ds)   [terminator]                         {Start};
%     \node (din)  [io,      right=2.2cm of ds]         {Input ciphertext, key};
%     \node (dkey) [process, right=2.2cm of din]        {Build key[],\; mod $\leftarrow$ length(key)};

%     \node (dinit) [process, below=1.7cm of ds]        {$i \leftarrow 0,\; len \leftarrow$ length(ciphertext)};
%     \node (dloop) [decision,below=1.7cm of din]       {$i < len$ ?};
%     \node (dout)  [io,      below=1.7cm of dkey]      {Output plaintext};

%     \node (dupq)  [decision,below=1.8cm of din,xshift=-2.2cm] {Uppercase?};
%     \node (dlowq) [decision,below=1.8cm of din]              {Lowercase?};
%     \node (dcopy) [process, below=1.8cm of din,xshift= 2.2cm] {plaintext[i] $\leftarrow$ ciphertext[i]};

%     \node (dup)   [process, below=1.8cm of dupq]      {plaintext[i] $\leftarrow$ ((c[i]-`A') - key[i \% mod] + 26) \% 26 + `A'};
%     \node (dlow)  [process, below=1.8cm of dlowq]     {plaintext[i] $\leftarrow$ ((c[i]-`a') - key[i \% mod] + 26) \% 26 + `a'};
%     \node (dinc)  [process, below=1.8cm of dcopy]     {$i \leftarrow i+1$};

%     \node (dret)  [process, below=1.8cm of dup]       {Return plaintext};
%     \node (dend)  [terminator, right=2.2cm of dret]   {End};
%     \node (dpad)  [process, right=2.2cm of dend]      {Decryption complete};

%     % Main flow
%     \draw [line] (ds) -- (din);
%     \draw [line] (din) -- (dkey);
%     \draw [line] (dkey) |- (dinit);
%     \draw [line] (dinit) -- (dloop);

%     % Loop decision
%     \draw [line] (dloop.east) -- node[above]{No} (dout.west);
%     \draw [line] (dout) |- (dret);

%     \draw [line] (dloop.south) -- node[right]{Yes} (dlowq.north);
%     \draw [line] (dlowq.west) -- node[above]{No} (dupq.east);
%     \draw [line] (dlowq.east) -- node[above]{No} (dcopy.west);
%     \draw [line] (dlowq.south) -- node[right]{Yes} (dlow.north);

%     \draw [line] (dupq.south) -- node[right]{Yes} (dup.north);

%     \draw [line] (dup.east) -- ++(0.9,0) |- (dinc.west);
%     \draw [line] (dlow.east) -- ++(0.9,0) |- (dinc.west);
%     \draw [line] (dcopy.south) -- (dinc.north);

%     % Loop back (routed outside, no node crossing)
%     \draw [line] (dinc.east) -- ++(1.2,0) |- (dloop.east);

%     % End path
%     \draw [line] (dret) -- (dend);
%     \draw [line] (dend) -- (dpad);
%   \end{autofitflowchart}
%   \caption{Vigenere Cipher Decryption Flowchart}
% \end{figure}

\section{Code}
\begin{minted}{cpp}
#include<bits/stdc++.h>
using namespace std;
class vigenere {
  private:
    string key_string;
    vector<int> key;
    int key_length;
  public:
    vigenere(string k);
    string encrypt(string s);
    string decrypt(string s);
};
vigenere::vigenere(string k) {
  string key_string = k;
  int ctr = 0;
  key.resize(k.length());
  for (int i = 0; i < k.length(); i++) {
    if (isupper(k[i])) {
      key[ctr] = k[i] - 'A';
      ctr++;
    }
    if (islower(k[i])) {
      key[ctr] = k[i] - 'a';
      ctr++;
    }
  }
  key_length = ctr;
}
string vigenere::encrypt(string s) {
  string ans = s;
  int len = s.length();
  int key_index = 0;
  for(int i = 0; i < len; i++) {
    if (s[i] >= 'a' && s[i] <= 'z') {
      ans[i] = ((s[i] - 'a') + key[key_index]) % 26 + 'a';
      key_index = (key_index + 1) % key_length;
    } else if (s[i] >= 'A' && s[i] <= 'Z') {
      ans[i] = ((s[i] - 'A') + key[key_index]) % 26 + 'A';
      key_index = (key_index + 1) % key_length;
    } else {
      ans[i] = s[i];
    }
  }
  return ans;
}
string vigenere::decrypt(string s) {
  string ans = s;
  int len = s.length();
  int key_index = 0;
  for(int i = 0; i < len; i++) {
    if (s[i] >= 'a' && s[i] <= 'z') {
      ans[i] = ((s[i] - 'a') - key[key_index] + 26) % 26 + 'a';
      key_index = (key_index + 1) % key_length;
    } else if (s[i] >= 'A' && s[i] <= 'Z') {
      ans[i] = ((s[i] - 'A') - key[key_index] + 26) % 26 + 'A';
      key_index = (key_index + 1) % key_length;
    } else {
      ans[i] = s[i];
    }
  }
  return ans;
}
int main() {
  cout << "Enter a string(with no blank spaces) key: ";
  string s;
  cin >> s;
  cin.ignore();
  cout << "Key is: " << s << endl;
  vigenere cipher(s);
  cout << "Enter Plaintext : ";
  string text;
  getline(cin, text);
  string encrypted = cipher.encrypt(text);
  cout << "Encrypted text: " << encrypted << endl;
  string decrypted = cipher.decrypt(encrypted);
  cout << "Decrypted text: " << decrypted << endl;
  return 0;
}
\end{minted}

\section{Result and Discussion - I/O}
The key is directly converted so no non-alphabetic characters are allowed.
\begin{minted}{text}
Enter a string(with no blank spaces) key: Thisisthekey
Key is: Thisisthekey
Enter Plaintext : Sadman Ishtiak
Encrypted text: Lhleif Bzldmyd
Decrypted text: Sadman Ishtiak
\end{minted}

\section{Conclusion}
The Vigen\`ere cipher improves on additive ciphers by using a repeating string key to vary the shift per character. It is still vulnerable to frequency analysis with enough ciphertext, but effectively demonstrates the concept of polyalphabetic encryption.
